\documentclass[11pt]{article}
\usepackage[a4paper,margin=1.9cm]{geometry}
\usepackage[T1]{fontenc}
\usepackage{textcomp}
\usepackage[spanish]{babel}
\usepackage{amsmath,amssymb}
\usepackage{lmodern}
\renewcommand{\familydefault}{\sfdefault}
\usepackage{enumitem}
\newcommand{\vect}[2]{\begin{pmatrix}#1\\#2\end{pmatrix}}
\newcommand{\vectiii}[3]{\begin{pmatrix}#1\\#2\\#3\end{pmatrix}}
\usepackage[table]{xcolor}
\usepackage{parskip}
\definecolor{unalblue}{RGB}{31,73,124}
\newlist{opciones}{enumerate}{1}
\setlist[opciones]{label=\alph*),itemsep=4pt,topsep=4pt,leftmargin=1.8em,font=\normalfont}
\setlist[enumerate,1]{itemsep=1em,topsep=0.8em}
\newcommand{\examfooter}{\vfill\rule{\linewidth}{0.4pt}\\[1pt]{\scriptsize\textit{Transcripción digital --- MathUNAL}\hfill\textcolor{unalblue}{\textbf{\thepage}}}}
\pagestyle{empty}

\begin{document}
\noindent
\begin{minipage}[t]{0.6\linewidth}
{\small\bfseries Geometría Vectorial y Analítica --- Parcial 3}\\
{\small Semestre 2022--02}
\end{minipage}%
\hfill
\begin{minipage}[t]{0.35\linewidth}
\raggedleft\small Escuela de Matemáticas. Universidad Nacional de Colombia, Sede Medellín.
\end{minipage}\\[2pt]
\rule{\linewidth}{0.6pt}

\begin{center}
\textbf{Geometría Vectorial --- Parcial 3}\\
Octubre de 2022
\end{center}

\textbf{Instrucciones:} La duración del examen es \textbf{1:50 horas}. Verifique que sus datos de identificación son correctos. La prueba consta de cinco preguntas. Escriba sus respuestas a los ejercicios 1,2,3,4 en la tabla de puntaje. No use fracciones, solamente decimales y redondee a dos cifras decimales. La pregunta 5 es de procedimiento, y su respuesta debe escribirse debajo del enunciado. No está permitido: hacer preguntas, usar calculadoras ni trabajar en hojas diferentes a las del examen. No se permite sacar el teléfono celular durante el examen.

\begin{center}
\begin{tabular}{|c|c|c|}
\hline
Ejercicio & Porcentaje & Respuesta\\
\hline
1 & 16\% & \\
\hline
2 & 16\% & \\
\hline
3 & 20\% & \\
\hline
4 & 18\% & $t_1=$\hspace{1.5cm}$t_2=$\\
\hline
5 & 30\% & Procedimiento\\
\hline
\end{tabular}
\end{center}

\begin{enumerate}
\item De las siguientes afirmaciones, seleccione todas aquellas que son verdaderas.
\begin{opciones}
\item Si $A=\vect{1}{1/3}$, $B=\vect{1/3}{-1}$, y $C=\vect{-1}{-1/3}$, entonces: $A$ y $B$ son linealmente independientes, $B$ y $C$ son linealmente independientes, pero $A$ y $C$ son linealmente dependientes.
\item Si $A=\vect{1}{1/3}$, $B=\vect{0}{0}$, y $C=\vect{1/3}{-1}$, entonces: $A$ y $B$ son linealmente dependientes, $B$ y $C$ son linealmente dependientes, y $A$ y $C$ son linealmente dependientes.
\item Sean $A$, $B$ y $C$ vectores en $\mathbb{R}^2$. Si $A$ y $B$ son linealmente dependientes, $B$ y $C$ son linealmente independientes y, $A$ y $C$ son linealmente dependientes entonces $A$ debe ser el vector nulo.
\item Si $A=\vect{1}{1/3}$, $B=\vect{1/3}{1/9}$, y $C=-11/3$, entonces: $A$ y $B$ son linealmente dependientes, $B$ y $C$ son linealmente independientes, y $A$ y $C$ son linealmente dependientes.
\item Ninguna de las anteriores.
\end{opciones}

\item Suponga que $A$ y $B$ son puntos en el plano y que $T$ es una transformación lineal. Si se sabe que $\det(T)=2$, y
$$T(A)=\vect{-2}{4},\qquad T(B)=\vect{4}{2},$$
¿cuál es el área del triángulo $\triangle AOB$?

\item Se sabe que una transformación lineal $T$ verifica las condiciones
$$T\vect{-1}{1}=\vect{0}{0},\qquad T\vect{2}{3}=\vect{5}{-5}.$$
Escoja todos los enunciados que son verdaderos.
\begin{opciones}
\item La transformación $T$ es invertible.
\item La imagen del segmento de recta con extremos $\vect{2}{3}$ y $\vect{3}{-1}$ es el punto $\vect{5}{-5}$.
\item El núcleo de la transformación $T$ es la línea recta con pendiente $3/2$.
\item El rango de la transformación $T$ es todo $\mathbb{R}^2$.
\item La imagen bajo la transformación lineal $T$ del triángulo con vértices $\vect{-3}{0}$, $\vect{0}{0}$, $\vect{0}{2}$ es un segmento.
\item Ninguna de las anteriores.
\end{opciones}

\item Sea $A=\vectiii{2}{-2}{t}$ y $\mathcal{P}$ el plano con ecuación general $-2x+(-2)y+(-1)z+(-1)=0$. Determine para qué valores de $t$ la distancia del punto $A$ al plano $\mathcal{P}$ es 1.

\item \textbf{Pregunta de procedimiento:} En el espacio de abajo escriba un procedimiento completo justificando cada paso. Considere los planos $\mathcal{P}_1$ y $\mathcal{P}_2$ con ecuaciones
$$\mathcal{P}_1: x+5y+3z-2=0;\qquad \mathcal{P}_2: X=\vectiii{4}{0}{0}+s\vectiii{-2}{1}{-1}+t\vectiii{3}{0}{-1},\ s,t\in\mathbb{R}.$$
\begin{enumerate}[label=\alph*)]
\item Verifique que $\mathcal{P}_1$ y $\mathcal{P}_2$ son paralelos.
\item Compruebe que $\mathcal{P}_1$ y $\mathcal{P}_2$ no se interceptan.
\item Halle una ecuación vectorial paramétrica de la recta $\mathcal{L}$ que es perpendicular a $\mathcal{P}_2$ y que pasa por el punto $P=\vectiii{4}{0}{0}$.
\end{enumerate}
\end{enumerate}

\examfooter
\end{document}
